\setlength{\headheight}{29.49297pt}

\section{Desarrollo}
En esta sección se presentan los detalles del proceso de modelización de la trayectoria de la nave Orión, junto con los criterios adoptados en cada etapa.

\subsection{Órbita Lunar}

Como paso previo a la simulación de la trayectoria de Orión, se modeló la órbita de la Luna alrededor de la Tierra. Esto permite conocer la posición lunar en todo instante, dato necesario para evaluar su influencia gravitatoria sobre la nave.

\subsubsection{Consigna}
La trayectoria de la nave se ve fuertemente afectada por el campo gravitatorio de la Tierra y de la Luna.
Para ello, se necesita saber la posición de la Luna en todo momento. Asumiendo posición inicial y
velocidad, se puede calcular la trayectoria con la siguiente ecuación diferencial:

\begin{equation}
    \frac{d^2\vec{x}}{dt^2} = \vec{a}
\end{equation}

Donde $\vec{x}$ es la posición vectorial de la Luna y $\vec{a}$ la aceleración producida por el campo
gravitatorio terrestre. Si escribimos la posición y la aceleración en dos dimensiones (x e y), y
descomponemos la ecuación de segundo orden en dos ecuaciones de primer orden, nos queda:

\begin{align}
    \frac{dx}{dt} &= v_x \\[6pt]
    \frac{dy}{dt} &= v_y \\[6pt]
    \frac{dv_x}{dt} &= G\,\frac{M_T}{d_T^2}\cos\alpha \\[6pt]
    \frac{dv_y}{dt} &= G\,\frac{M_T}{d_T^2}\operatorname{sen}\alpha
\end{align}

Con $G$ la constante de gravitación universal, $M_T$ la masa de la Tierra, $d_T$ la distancia
entre la Luna y la Tierra, y $\alpha$ el ángulo que forma el vector posición con el eje $x$.

\subsubsection{Discretización de EDOs}
Para resolver las ecuaciones dadas por métodos numéricos mediante PVI se discretizaron las ecuaciones de la siguiente forma:

\paragraph{Euler Explícito}
Dado el sistema de ecuaciones de primer orden, la discretización por el método de Euler explícito resulta:

\begin{align}
    x_{n+1}      &= x_n + h\, v_{x,n} \\[4pt]
    y_{n+1}      &= y_n + h\, v_{y,n} \\[4pt]
    v_{x,n+1}    &= v_{x,n} + h\, G\,\frac{M_T}{d_{T,n}^2}\cos\alpha_n \\[4pt]
    v_{y,n+1}    &= v_{y,n} + h\, G\,\frac{M_T}{d_{T,n}^2}\operatorname{sen}\alpha_n
\end{align}

donde $h$ es el paso de integración y $n$ el índice del paso temporal.

\paragraph{Runge-Kutta de orden 2 (RK2)}
Definiendo $\vec{u} = (x,\, y,\, v_x,\, v_y)^T$ y $\vec{f}(\vec{u})$ como el lado derecho del sistema, las etapas del método son:

\begin{align}
    \vec{q}_1 &= \vec{f}\!\left(\vec{u}_n\right) \\[4pt]
    \vec{q}_2 &= \vec{f}\!\left(\vec{u}_n + h\,\vec{q}_1\right) \\[6pt]
    \vec{u}_{n+1} &= \vec{u}_n + \frac{h}{2}\left(\vec{q}_1 + \vec{q}_2\right)
\end{align}

\paragraph{Runge-Kutta de orden 4 (RK4)}
Las cuatro etapas del método clásico de Runge-Kutta son:

\begin{align}
    \vec{q}_1 &= \vec{f}\!\left(\vec{u}_n\right) \\[4pt]
    \vec{q}_2 &= \vec{f}\!\left(\vec{u}_n + \frac{h}{2}\,\vec{q}_1\right) \\[4pt]
    \vec{q}_3 &= \vec{f}\!\left(\vec{u}_n + \frac{h}{2}\,\vec{q}_2\right) \\[4pt]
    \vec{q}_4 &= \vec{f}\!\left(\vec{u}_n + h\,\vec{q}_3\right) \\[6pt]
    \vec{u}_{n+1} &= \vec{u}_n + \frac{h}{6}\left(\vec{q}_1 + 2\vec{q}_2 + 2\vec{q}_3 + \vec{q}_4\right)
\end{align}

\subsubsection{Condiciones iniciales del PVI}

Se coloca a la Luna en el \textbf{perigeo} como condición inicial, es decir, en el punto de mínima
distancia a la Tierra. En ese instante el vector posición es puramente horizontal y la velocidad
es perpendicular al radio:

\[
x_0 = R_{\text{perigeo}}, \quad y_0 = 0, \quad v_{x,0} = 0, \quad v_{y,0} = v_{\text{perigeo}}
\]

Los valores de perigeo y apogeo se obtienen de datos observacionales de la órbita lunar
media \cite{nasa_horizons}:

\[
R_{\text{perigeo}} = 356\,500 \text{ km} = 3{,}565 \times 10^{8} \text{ m}, \qquad
R_{\text{apogeo}}  = 406\,700 \text{ km} = 4{,}067 \times 10^{8} \text{ m}
\]

La velocidad en el perigeo se obtiene de la \textbf{ecuación de la vis-viva}:

\[
v = \sqrt{G M_T \left(\frac{2}{r} - \frac{1}{a}\right)}
\]

donde $a$ es el semieje mayor de la órbita elíptica:

\[
a = \frac{R_{\text{perigeo}} + R_{\text{apogeo}}}{2} = 3{,}816 \times 10^{8} \text{ m}
\]

Evaluando en $r = R_{\text{perigeo}}$ con las constantes
$G = 6{,}674 \times 10^{-11}\ \text{m}^3\,\text{kg}^{-1}\,\text{s}^{-2}$ y
$M_T = 5{,}972 \times 10^{24}\ \text{kg}$:

\[
v_{\text{perigeo}} = \sqrt{G M_T \left(\frac{2}{R_{\text{perigeo}}} - \frac{1}{a}\right)}
\approx 1{,}076 \times 10^{3}\ \text{m/s}
\]

Resumiendo, el vector de estado inicial es:

\[
\boxed{
\mathbf{q}(0) = \begin{pmatrix} 3{,}565 \times 10^{8}\ \text{m} \\ 0 \\ 0 \\ 1{,}076 \times 10^{3}\ \text{m/s} \end{pmatrix}
}
\]

\subsubsection{Elección del paso temporal}

Para obtener la trayectoria de una vuelta completa de la Luna en torno a la
Tierra, se integró durante un período orbital sidéreo de $T = 27{,}32$ días
$\approx 2{,}360{,}045$s. El paso de tiempo $h$ se
eligió de forma que el número de pasos $N = T/h$ sea suficientemente grande
para que el error de truncamiento local no acumule un error global apreciable.
Se adoptó $h = 60$s, que representa una fracción $h/T \approx 2{,}5 \times 10^{-5}$
del período orbital, garantizando al menos $N \approx 39{,}000$ puntos por
órbita. Este nivel de discretización es suficiente para resolver con precisión
las variaciones de posición y velocidad a lo largo de la trayectoria.

\subsubsection{Resultados}
Se llegó a los siguientes resultados:
